\section{The ordinary theory and the iPRA interpretations}
\label{sec:interpretations}

\subsection{Ordinary realization of the complete calculus}
The ordinary reading of the paired binder is
\[
 (D_{x,y}B)^\circ
 =\exists x\bigl(B^\circ(x,0)\to\forall y\,B^\circ(x,y)\bigr),
\]
and $[A;y.B]^\circ=A^\circ\to\forall yB^\circ(y)$. Equality reads as equality. Each fixed recursor description has a first-order primitive-recursive realization. A bound update variable is a finite-syntax device, not a variable ranging over functions. No universal evaluator for primitive-recursive descriptions is postulated.

R is sound by its defining recursion equation. T is ordinary equality transport followed by symmetry. U is implication elimination followed by universal elimination. E is existential elimination with the stated fresh parameter. In I's case mode, numerical coverage extends the zero and successor cases to every number. In its equation-progress mode, ordinary quantifier-free induction, under the conditional zero instance, gives the required universal equation. Existential introduction supplies the displayed witness. This establishes direct soundness in iPRA, including the nontrivial side conditions.

The normalizing equation
\[
 R(0,a;u.h)\lit R(0,a;u.0)
\]
is inherited from the canonical representation. In ordinary arithmetic it says that both expressions are $Sa$. It must be implemented during substitution when a counter becomes literal zero; it has not disappeared from the trusted theory.

\subsection{Explicit numerical maps}
For each fixed $R$ description, the forward map to ordinary primitive-recursive notation is
\[
 R(n,a;u.h)\longmapsto I(n,Sa;u.h^\circ),
\]
where ordinary initialized iteration satisfies $I(0,a;h)=a$ and $I(Sn,a;h)=h(I(n,a;h))$. Variables, zero, and equality are unchanged.

For the other direction, the inherited successor-seeded bootstrap constructs one fixed finite term $P$ and proves internally
\[
 P0=0,\qquad P(Sa)=a.
\]
Its complete construction and proof are preserved in \file{reference/successor_seed_bootstrap.tex}; its finite-case prerequisites are in \file{reference/successor_seed_prelogic.tex}. These are results of the current proposal, reused here, not new assumptions. Define
\[
 \widehat I(n,a;u.h)
 :=P\!\left(R\bigl(n,a;v.S(h[P(v)/u])\bigr)\right).
\]
The zero computation is $P(Sa)=a$. The successor computation follows by R and then $P(Sw)=w$. Consequently this single fixed translation recovers initialized scalar iteration, including its independent initial value.

The inherited scalar development supplies a primitive-recursive pairing $\pi$ with projections $L,H_1$ satisfying
\[
 L(\pi(a,b))=a,\qquad H_1(\pi(a,b))=b.
\]
These are numerical terms on the same domain; the domain is not replaced by their coding range. For a fixed primitive-recursive definition
\[
 f(0,\vec p)=g(\vec p),\quad
 f(Sn,\vec p)=h(n,f(n,\vec p),\vec p),
\]
its translated term is
\[
 \begin{split}
 W(v,\vec p)&=\pi(SL(v),h^*(L(v),H_1(v),\vec p)),\\
 V(n,\vec p)&=\widehat I(n,\pi(0,g^*(\vec p));v.W(v,\vec p)),\\
 f^*(n,\vec p)&=H_1(V(n,\vec p)).
 \end{split}
\]
The zero, successor, projection and composition descriptions translate directly. This is an external recursion on each finite function description. The proof of its defining equations uses internal equation induction with all passive parameters fixed.

\subsection{Full logical expressiveness}
Let $C(n,a,b)=\widehat I(n,a;u.b)$. Then $C(0,a,b)=a$ and $C(Sn,a,b)=b$. The inherited clause construction~\cite{clauses} defines a formula gate $G_vA$ and formula selector $M_v(A,B)$, all as finite Eq/$D$ formulas, with
\[
 \vdash G_0A,\qquad G_{Sv}A\dashv\vdash A,\qquad
 M_0(A,B)\dashv\vdash A,\qquad M_{Sv}(A,B)\dashv\vdash B.
\]
For equations, gate each side with $C(v,0,-)$ and select their two sides with $C$. For two $D$ formulas, recurse through aligned fresh binders. If just one selected formula is a $D$ formula, pad the equation with its gate under an unused $D$ before aligning the binders. Thus the recursion terminates on the maximum binder depth. The padding equivalence is $A\dashv\vdash D_{w,y}G_yA$.

Put $p(z)=\widehat I(z,0;u.C(u,1,0))$, so $p(1)=1$ and $p(2)=0$. The full logical macros are
\[
\begin{aligned}
 A\Rightarrow_*B&=D_{w,z}M_z(A,B),\\
 A\wedge_*B&=D_{w,z}G_z M_{p(z)}(A,B),\\
 \exists_*x\,A(x)&=D_{x,y}G_yA(x),\\
 A\vee_*B&=\exists_*x\,M_x(A,B),\\
 \forall_*x\,A(x)&=A(0)\wedge_*D_{w,z}G_zA(z),\\
 \bot_*&=(0=S0).
\end{aligned}
\]
All newly bound variables are fresh and $w$ is unused. These are abbreviations, not extra primitives in the verifier. Their usual intuitionistic introduction and elimination rules follow from the derived case rule, clause use, and paired elimination. For example, implication introduction is case introduction using the zero selector as its assumption and the successor selector as its conclusion; implication elimination opens its clause and specializes at $1$. Existential elimination specializes the opened gated clause at $1$, using the theorem at its zero gate. Conjunction projections use the two instances $1$ and $2$.

The native logical derivations, arithmetic equality reflection, and quantifier-free induction are included in the frozen references. They are unchanged because the new calculus derives exactly the old Eq/$D$ consequences. Progress remains restricted to a literal equation: allowing it for every compound matrix would introduce full induction and would not be this presentation of iPRA.

\subsection{Why these are bi-interpretations}
There are two separate round-trip obligations. The new logical translation removes $Q$ by $Q_yB\mapsto D_{w,y}T(B)$; Section~\ref{sec:clause-five} proves its round trips internally. On numerical syntax it is the identity.

For the inherited arithmetic maps, compare each fixed operation with its double translation in the external order of finite descriptions. Once earlier operations have been identified, the two trajectories have the same initial value and provably equal updates. Literal-equation induction proves their equality at every internal counter. This also justifies changes of a recursor body: they follow from a proved trajectory comparison, not an uncharged function-extensionality axiom. The same argument covers ordinary primitive recursion after compilation by pairing.

Formula structural induction then lifts the numerical identities and logical macro equivalences to every formula. Hence the numerical relation $x=y$ is internally a total bijection preserving all interpreted operations in each composite. This is stronger than merely knowing that two theories have the same theorems about standard integers, the same consistency strength, or mutual interpretations. Because falsity is preserved and the finite proof translations run both ways, equiconsistency follows as well.
